% Nouveau EGA V (cote 25): archive cote 25, site page 47 (PDF page 48).
% typescript. Draft compared: none.
% AI-transcribed by Claude Opus 5.5 (claude-opus-5-5, adaptive thinking, effort high), Message Batches API, batch ega25b (msgbatch_01A7n5YuXubUPNK1zsepi8RG), 2026-11-13; not yet checked by eye.
\scanpage{25/47}
\var{\margin{(\unsure{$\uparrow$|?} la variété \unsure{lin} L de codim 2 des hyperplans tangents à X en x, donc par raisons de dim. on a T = L, donc $X \subset L^{\circ}$ = polaire de L, qui est une droite, donc X est un \unsure{\uline{ouvert}} d'une droite.)}}{\placedmargin{(\unsure{$\uparrow$|?} la variété \unsure{lin} L de codim 2 des hyperplans tangents à X en x, donc par raisons de dim. on a T = L, donc $X \subset L^{\circ}$ = polaire de L, qui est une droite, donc X est un \unsure{\uline{ouvert}} d'une droite.)}}{marginal note placed}{ega25_25_047-v0001}
\mspage{48}
\unsure{\hand{(}}En langage classique, elle s'exprime \var{\xout{xxxdisant}}{}{typewriter strike-out}{ega25_25_047-v0002} sauf erreur en disant
que X est \uline{réglée} pour l'immersion projective considérée \hand{[}- et si ça nous
chante, nous avons ici  tout ce qu'il faut, grâce à 8.5. et ses corollai
res, pour expliquer et justitffier une telle terminologie, au cas où cela
t'inspirerait de faire le lien avec la taupe\var{\over{)}{]}}{]}{writing over}{ega25_25_047-v0003}. \var{\inserted{Par exemple, si dim $X=1$, cela implique que X est une droite, car si $x \in X(k)$, alors T contient \unsure{$\uparrow$|P}}}{Par exemple, si dim $X=1$, cela implique que X est une droite, car si $x \in X(k)$, alors T contient \unsure{$\uparrow$|P}}{insertion}{ega25_25_047-v0004}

b)  Lorsque la car. est p >0, il faudrait pour bien faire donner
des exemples (avec dim P = 2, X une courbe non singulière) où  les
conditions de 8.6. ne sont pas vérifiées i.e. $\underline{Y}^{\mathrm{sing}} = \underline{Y}^{\mathrm{sup\ sing}}$, et où
néanmoins  dim T = r-1 , i.e. \var{\inserted{des exemples}}{des exemples}{insertion}{ega25_25_047-v0005} où  L/K est une extension finie
inséparable.  J'ai la flemme de construire des exemples, mais je ne
doute guère qu'il y en ait.

\var{\deleted{\uline{Corollaire} 8.11.}}{}{deletion}{ega25_25_047-v0006} Dans a), faire un appel du pied au \var{\xout{xxxxx}}{}{typewriter strike-out}{ega25_25_047-v0007} nº suivant,
où on montrera que lorsque le cas exceptionnel ``réglé'' se produit, alors \var{\margin{N.B. X est dit ``réglé'' pour f si $\mathrm{Im}(\underline{Y}^{\mathrm{sing}} \to \check{P})$ est de codim $\unsure{\geq} 2$}}{\placedmargin{N.B. X est dit ``réglé'' pour f si $\mathrm{Im}(\underline{Y}^{\mathrm{sing}} \to \check{P})$ est de codim $\unsure{\geq} 2$}}{marginal note placed}{ega25_25_047-v0008}
par une modification triviale de l'immersion projective considérée, on
se retrouve dans la situation  ``générale''  de 8.8.

\var{\margin{\unsure{$\|$}}}{\placedmargin{\unsure{$\|$}}}{marginal note placed}{ega25_25_047-v0009}La partie du présent nº de 8.\var{\over{\illegible{1}}{6}}{6}{writing over}{ega25_25_047-v0010}. à ici  devrait sans doute faire
un nº à part, de nature différentielle,  tandis que le début du nº avec
ce qui va suivre serait bloqué dans un nº  sur  la dimension \unsure{dus|des}
exceptionnels. \hand{(Je n'utilise que le fait que $\underline{Y}^{\mathrm{sing}}$ est de
dim r-1).}

\uline{Proposition} 8.1\var{\over{2}{3}}{3}{writing over}{ega25_25_047-v0011}.   Nous supposons tou\var{\xout{x}}{}{typewriter strike-out}{ega25_25_047-v0012}jours f:X$\to$P non ramifié, \var{\inserted{et X sans pt isolé,}}{et X sans pt isolé,}{insertion}{ega25_25_047-v0013} \var{\deleted{Suppo-}}{}{deletion}{ega25_25_047-v0014}
\var{\inserted{On suppose que X satisfait $(R_k)$ géométrique\var{\deleted{, et si}}{}{deletion}{ega25_25_047-v0016}}}{On suppose que X satisfait $(R_k)$ géométrique\var{\deleted{, et si}}{}{deletion}{ega25_25_047-v0016}}{insertion}{ega25_25_047-v0015}
\var{\deleted{sons que X satisfasse la propriété $(R_k)$ géométrique  (où k \unsure{$\geq$} 0)\hand{,} et si}}{}{deletion}{ega25_25_047-v0017}
\var{\deleted{X' est \var{\xout{xxensemble}}{}{typewriter strike-out}{ega25_25_047-v0019} l'ouvert de X des points où X est lisse sur k, \unsure{le}}}{}{deletion}{ega25_25_047-v0018}
\var{\deleted{\var{\xout{\illegible{12}}}{}{typewriter strike-out}{ega25_25_047-v0021}}}{}{deletion}{ega25_25_047-v0020}
\var{\deleted{\var{\inserted{on suppose que le}}{on suppose que le}{insertion}{ega25_25_047-v0023}morphisme  induit f':X'  P  \var{\xout{soit}}{}{typewriter strike-out}{ega25_25_047-v0024} satisfasse \var{\xout{auxxconditionsxéquivalentes}}{}{typewriter strike-out}{ega25_25_047-v0025}}}{}{deletion}{ega25_25_047-v0022}
\var{\deleted{\var{\xout{\illegible{10}}}{}{typewriter strike-out}{ega25_25_047-v0027} k(  )}}{}{deletion}{ega25_25_047-v0026}
\var{\deleted{\var{\xout{dex8x8xxx(disonsxdoncxsimplementxX'xestxnonxx}}{}{typewriter strike-out}{ega25_25_047-v0029}  à la condition}}{}{deletion}{ega25_25_047-v0028}
\var{\deleted{g( $\underline{Y}'^{\,\mathrm{\unsure{sup\,sing}}}$) \hand{\unsure{´}}de codimension $\geq$2 dans $\hand{\check{P}}$ (condition \uline{automatiquement}}}{}{deletion}{ega25_25_047-v0030}
\var{\deleted{\var{\xout{v}}{}{typewriter strike-out}{ega25_25_047-v0032}vérifiée si k est de car 0 en vertu de de 8.11. ).}}{}{deletion}{ega25_25_047-v0031}  Soit $Z_k$ la parti
de \hand{$\check{P}$} \var{\xout{fermédesxxxxxx}}{}{typewriter strike-out}{ega25_25_047-v0033} complémentaire de l'ensemble des \hand{$\xi$}$\in$\hand{$\check{P}$}  tels

\var{\margin{47}}{\placedmargin{47}}{marginal note placed}{ega25_25_047-v0034}
